\documentclass{article}\usepackage[utf8]{inputenc}\usepackage[T1]{fontenc}\usepackage[french]{babel}\usepackage{geometry}\geometry{margin=2.5cm}\begin{document}\section*{Diagnostic et préconisations de sécurité - MSAP}\subsection*{1. Configurations à risque (Documents 1 et 2)}
L'examen des configurations actuelles révèle des failles de sécurité majeures :
\begin{itemize}
\item \textbf{Comptes locaux résiduels} : Bien que les utilisateurs s'identifient via l'Active Directory , des comptes locaux tels que "Administrateur", "DefaultAccount" et "Invité" sont présents sur les machines. Le compte "Invité" actif constitue une vulnérabilité critique.
\item \textbf{Élévation de privilèges par héritage} : Le groupe "partenaires" est membre du groupe "Administrateurs". Par conséquent, tous les membres des groupes partenaires (clic, enedis, msa, tresor) héritent de droits d'administration totale sur le domaine.
\item \textbf{Absence de cloisonnement} : Tous les partenaires sont regroupés dans une structure commune ("partenaires"), ce qui favorise les accès transversaux non autorisés.
\end{itemize}\subsection*{2. Synthèse des bonnes pratiques (Documents 3 et 4)}
Pour corriger ces failles, les mesures suivantes doivent être appliquées :
\begin{itemize}
\item \textbf{Principe du moindre privilège} : Les utilisateurs ne doivent disposer que des accès nécessaires à leurs missions.
\item \textbf{Restriction des accès partagés} : Les autorisations de partage pour le groupe "Tout le monde" (Everyone) doivent être supprimées. Actuellement, ce groupe dispose de permissions sur le dossier "PartageTelecentre".
\item \textbf{Gestion granulaire NTFS} : Il est impératif de définir des droits spécifiques pour chaque dossier partenaire (ex: seul le groupe ENEDIS accède au dossier \texttt{enedis}) au lieu d'accorder des droits globaux sur la racine du partage.
\end{itemize}\subsection*{3. Risques induits par les privilèges excessifs (Document 4)}
L'attribution de droits trop larges (Modfication/Contrôle total) expose la MSAP à :
\begin{itemize}
\item \textbf{Compromission de l'intégrité} : Un utilisateur peut modifier ou supprimer des données critiques appartenant à d'autres partenaires.
\item \textbf{Risque d'attaque latérale} : Si un compte du groupe "partenaires" est compromis, l'attaquant dispose immédiatement de privilèges administrateur pour s'étendre sur le réseau.
\end{itemize}\subsection*{4. Préconisations pour la segmentation (Document 5)}
La sécurisation réseau doit s'appuyer sur une isolation logique et filtrée :
\begin{itemize}
\item \textbf{Segmentation par VLAN} : Maintenir l'isolation des partenaires dans des VLAN distincts ($VLAN 10, 20, 30, 40$).
\item \textbf{Filtrage par Access Control Lists (ACL)} : Configurer le routeur principal pour interdire les flux entre les VLANs des partenaires tout en autorisant l'accès au VLAN 100 du serveur de fichiers.
\item \textbf{Syntaxe de filtrage (Exemple Cisco)} : Pour bloquer un flux, utiliser la commande :
$$deny \ 192.168.X.0 \ 0.0.0.255 \ 192.168.Y.0 \ 0.0.0.255$$où $X$ et $Y$ sont les sous-réseaux des différents partenaires.
\end{itemize}\end{document}